\documentclass[10pt,twocolumn,a4paper]{article}

% --- Page geometry (CIM 2026: A4, 17.2x25.2cm text, 2cm top, 2.5cm bottom, 1.9cm sides) ---
\usepackage[a4paper,
  top=2.0cm, bottom=2.5cm,
  left=1.9cm, right=1.9cm,
  columnsep=0.8cm]{geometry}

\usepackage[T1]{fontenc}
\usepackage[utf8]{inputenc}
\usepackage[italian]{babel}
\usepackage{times}
\usepackage{graphicx}
\usepackage{amsmath}
\usepackage{url}
\usepackage{hyperref}
\usepackage{natbib}
\usepackage{caption}
\usepackage{float}
\usepackage{tikz}
\usetikzlibrary{positioning,arrows.meta,shapes.geometric,fit,calc}
\usepackage{listings}
\lstset{
  basicstyle=\ttfamily\fontsize{8}{9.5}\selectfont,
  columns=fullflexible,
  keepspaces=true,
  showstringspaces=false,
  aboveskip=4pt, belowskip=4pt,
  xleftmargin=4pt
}

% Font sizes per CIM 2026 template
\renewcommand{\normalsize}{\fontsize{10}{12}\selectfont}
\renewcommand{\small}{\fontsize{8}{10}\selectfont}

% No headers/footers/page numbers (added by proceedings editor)
\pagestyle{empty}

% Section heading styles
\usepackage{titlesec}
\titleformat{\section}{\normalfont\fontsize{12}{14}\bfseries\centering}{}{0em}{}
\titleformat{\subsection}{\normalfont\fontsize{10}{12}\bfseries}{}{0em}{}
\titleformat{\subsubsection}{\normalfont\fontsize{10}{12}\itshape}{}{0em}{}
\titlespacing{\section}{0pt}{12pt plus 2pt minus 2pt}{6pt plus 1pt}
\titlespacing{\subsection}{0pt}{6pt plus 2pt minus 2pt}{3pt plus 1pt}

% Caption style
\captionsetup{font=small, labelfont=bf, justification=justified}

% --- Title block ---
\makeatletter
\renewcommand{\maketitle}{%
  \twocolumn[{%
    \begin{center}
      {\fontsize{16}{18}\bfseries\MakeUppercase{\@title}\par}
      \vspace{1em}
      {\@author\par}
    \end{center}
    \vspace{1em}
  }]%
}
\makeatother

% --- Copyright footnote (first page, bottom left) ---
\usepackage{footnote}
\makeatletter
\def\blfootnote{\xdef\@thefnmark{}\@footnotetext}
\makeatother

% ================================================================
% NOTA: titolo + autore da finalizzare in fase di submission.
% Versione double-blind: titolo neutro, autore "Anonymous".
\title{Un ambiente di sintesi granulare in tempo differito:\\
       il loop lungo come postura compositiva}

\author{
  \begin{tabular}{c}
    Anonymous\\
    \small [Anonymous]\\
    \small anonymous@anonymous
  \end{tabular}
}
% ================================================================

\begin{document}
\maketitle

\blfootnote{\small \textit{Copyright: \copyright 2026 Giulio De Mattia Romano. This is an open-access article distributed under the terms of the \href{https://creativecommons.org/licenses/by/4.0/}{Creative Commons Attribution License 4.0 International}, which permits unrestricted use, distribution, and reproduction in any medium, provided the original author and source are credited.}}

% ----------------------------------------------------------------
\begin{abstract}
% 150--200 parole. Inglese. Da scrivere in fase finale, dopo che
% le sezioni 1--6 sono stabili.
\end{abstract}

% ================================================================
\section{Introduzione}

La sintesi granulare ha attraversato, nei suoi primi quindici anni di
implementazione computer, una migrazione netta dal tempo differito al
tempo reale. La prima implementazione documentata~\cite{Roads1978}
opera in modalità offline su B6700 ALGOL come front-end per MUSIC V:
il computer non regge la generazione real-time delle migliaia di grani
per secondo che il paradigma richiede, e l'unico modo possibile è
scrivere, compilare, ascoltare dopo. Nel contesto italiano la stessa
necessità è formulata esplicitamente da Di~Scipio nel 1991, a chiusura
della fase offline su microcomputer: \emph{«queste procedure sono
attualmente implementate in tempo differito, su un IBM PC 286 [\dots]
un problema attualmente insormontabile sta nella quantità di RAM nella
quale il segnale da ridurre in grani è conservato»}~\cite{DiScipio1991cim}.
Il tempo differito è, in questa fase, vincolo hardware, non scelta
compositiva.

Tre anni prima Truax~\cite{Truax1988} aveva descritto su DMX-1000 la
prima implementazione real-time stabile della granulazione. Il vincolo
hardware cade, e con esso, secondo l'autore, una postura compositiva:
\emph{«the key is to abandon linear modes of compositional thinking,
which result in deterministic output (e.g., score or sequencer driven),
and to substitute process-oriented multitask strategies for real-time
execution»} (p.~19). Il real-time è teorizzato non come miglioria
tecnica ma come cambio di paradigma: feedback immediato, abbandono
della partitura pre-scritta, processi multi-task.

Il sistema descritto in questo lavoro compie il percorso inverso: torna
volontariamente al tempo differito in un momento in cui il real-time è
disponibile. Il ritorno non è un passo indietro né una critica al
programma di Truax. Corrisponde a una postura compositiva personale e
situata: quella in cui composizione e studio della tecnica coincidono,
e in cui il \emph{loop di feedback lungo} --- specifica, generazione,
ascolto, riflessione, riscrittura --- è lo spazio necessario per
abitare gli spazi compositivi della granulazione come forma e
struttura. Il giudizio drammaturgico del compositore garantisce la
non-linearità del processo indipendentemente dalla scala temporale del
feedback: nel real-time opera sull'atto immediato, nel tempo differito
opera sulla riflessione tra un ciclo e l'altro. Cambia la scala, non la
natura. Il contributo specifico del sistema --- DSL YAML come
linguaggio di intenzioni parametriche, partitura grafica automatica
con asse Y centrato sulla posizione nel buffer sorgente, workflow STEMS
con cache incrementale --- è la realizzazione tecnica di questa
postura, ed eredita la prospettiva non neutrale degli strumenti
compositivi enunciata in ambito CIM da Arcella e Silvestri:
\emph{«tools and technologies used to produce a musical work are not
neutral but incorporate knowledge that influence the choices of the
composer»}~\cite{Arcella2012}.

% ================================================================
\section{Sintesi granulare: dal paradigma Gabor al controllo gerarchico}

\subsection{Il quanto acustico e la prima implementazione}

Il paradigma granulare si fonda sulla proposta di
Gabor~\cite{Gabor1947} di rappresentare il suono come matrice di quanti
elementari su un piano tempo--frequenza, soggetti al principio di
indeterminazione $\Delta t \cdot \Delta f \geq 1$. Ogni grano è
prodotto di una sinusoide e di una finestra di durata finita; un suono
complesso è ricostruito come collezione discreta di grani.

L'implementazione computer della granulazione introduce un problema
specifico: la specifica esplicita dei parametri di ogni singolo grano
diventa intrattabile non appena la densità supera poche unità per
secondo. Roads~\cite{Roads1978} lo affronta organizzando il sistema in
due livelli: un \emph{front-end} ad alto livello (AGS, scritto in
ALGOL) accetta una specifica compatta a sei coppie valore/slope per
ogni \emph{event} e genera la nota-list completa per MUSIC~V come
\emph{engine} di sintesi. Il pattern \emph{front-end~$\to$~engine},
che riemerge in ogni implementazione successiva di granulazione
algoritmica offline, è qui formulato per la prima volta in modo
esplicito.

\subsection{Il controllo gerarchico in tempo reale}

Truax~\cite{Truax1988} consolida la pratica granulare in tempo reale
sul DMX-1000 e ne fissa il vocabolario di controllo. Due elementi
sono centrali per il presente lavoro. Il primo è la gerarchia di
controllo (Fig.~3 nell'originale), che separa i parametri di sintesi
(operanti per-grano) dai parametri di organizzazione macro-formale
(operanti su finestre temporali estese): la specifica non avviene
grano per grano ma per regioni di parametro entro cui il singolo grano
è campionato. Il secondo è la \emph{Tabella~1 dei psychoacoustic
correlates}: ogni parametro di controllo (frequenza centrale, durata
del grano, densità, ampiezza, scatter) è associato a un corrispettivo
percettivo che è più significativo del valore numerico grezzo come
base per l'organizzazione compositiva.

La conseguenza estetica è formalizzata in Truax~\cite{Truax1994}: la
granulazione \emph{«links frequency and time at the micro level [so
it] makes it possible to treat these two variables independently at
the macro level»}~(p.~44). La separazione micro/macro non è
convenzione descrittiva ma proprietà psicoacustica abilitante, e
giustifica successivamente l'esperienza del time-stretching granulare
come spostamento dell'attenzione \emph{«inside the
sound»}~\cite{Truax2014}. Già nel 1990 Truax aveva inoltre dichiarato
inadeguata, anche fuori dal vincolo real-time, la specifica grano per
grano: \emph{«it would be absurd to specify the parameters for each
grain individually»}~\cite{Truax1990}, formulando il problema di
controllo che ogni linguaggio di specifica granulare deve risolvere.

\subsection{Il ramo CIM offline}

Parallelamente alla tradizione anglofona, il contesto italiano
sviluppa una linea autonoma di ricerca granulare offline. De~Poli e
Piccialli~\cite{DePoliPiccialli1988,DePoliPiccialli1991} elaborano un
modello di sintesi granulare additiva pitch-synchronous con grani
trattati come risposte FIR e controllo formantico dell'inviluppo
spettrale. Ortosecco e Piccialli~\cite{OrtoseccoPiccialli1989}
proseguono lo stesso filone identificando \emph{wavelet = grano} come
base teorica rigorosa e implementando un sistema di analisi/risintesi
con wavelet prototipo tabulata una sola volta su 4096 campioni e
grani ottenuti per sottocampionamento. Il pattern
\emph{precompute-once~/~reuse-many} è la realizzazione CIM 1989
diretta del componente che il presente sistema chiama
\texttt{WindowGenerator}.

Roads, nel primo paper CIM dedicato alla granulazione~\cite{Roads1985cim},
formalizza esplicitamente il problema d$\cdot$n (densità per durata) come
motivazione del DSL parametrico, introduce il \emph{frame} come unità
temporale sopra il grano e propone il \emph{polygon} su piano
frequenza/tempo come notazione per la traiettoria di un evento
granulare. Tutti e tre i concetti sono precursori diretti, sul piano
CIM, di costrutti del sistema descritto in \S\ref{sec:architettura}.

Di~Scipio~\cite{DiScipio1991cim}, oltre alla formulazione del vincolo
hardware già citata, sceglie come modello di controllo le mappe
caotiche deterministiche ($x_{n+1} = f(x_n)$, dipendenza esplicita del
grano successivo dal precedente). Si distingue per famiglia di
controllo dalla tradizione tendency-mask di Truax: \emph{caos iterativo}
contro \emph{distribuzione statistica campionata grano per grano}. Le
due famiglie convivono nella tradizione CIM offline come alternative
parallele, non come varianti dello stesso pattern; il sistema descritto
in questo lavoro eredita esplicitamente la seconda. Arcella e
Silvestri~\cite{Arcella2012}, ricostruendo \emph{Analogique~B} di
Xenakis, fissano in ambito CIM la topologia
\emph{algoritmo~$\to$~score~$\to$~Csound~$\to$~audio} esplicitamente
fattorizzata in due moduli (\emph{«our software implementation factors
the whole problem in two»}, p.~147) --- precursore più diretto della
pipeline che la \S\ref{sec:architettura} descrive, da cui il sistema
proposto si distingue per il livello di astrazione del modulo di
specifica (YAML dichiarativo + IR Python, non emissione diretta di
score Csound testuale).

% ================================================================
\section{Architettura per l'indagine parametrica}
\label{sec:architettura}

L'architettura del sistema descritto in questo lavoro è progettata
attorno al loop lungo: ogni componente è la realizzazione tecnica di
un anello specifico del ciclo specifica $\to$ generazione $\to$
ascolto $\to$ riflessione $\to$ riscrittura. La pipeline (Fig.~1)
separa una specifica testuale dichiarativa --- un file YAML --- da
una Intermediate Representation (IR) Python composta da oggetti
\texttt{Stream}, \texttt{Voice}, \texttt{Grain}, \texttt{Parameter},
e da un layer di rendering disaccoppiato. La topologia
\emph{specifica $\to$ IR $\to$ engine $\to$ audio} è ereditata dal
precursore CIM più diretto~\cite{Arcella2012}, che enuncia
esplicitamente la fattorizzazione del problema in due moduli
(\emph{«our software implementation factors the whole problem in
two»}, p.~147); rispetto a quella implementazione il sistema
proposto si distingue per il livello di astrazione del modulo di
specifica, che non emette score Csound testuale ma una IR
intermedia parametrizzata, manipolabile prima del rendering.

\begin{figure}[t]
  \centering
  \begin{tikzpicture}[
      node distance=4pt and 4pt,
      every node/.style={font=\fontsize{7}{8}\selectfont},
      box/.style={draw, rectangle, minimum height=4mm, inner sep=2pt,
                  align=center, font=\fontsize{7}{8}\ttfamily},
      lbl/.style={align=center, font=\fontsize{7}{8}\selectfont},
      arr/.style={-{Latex[length=1.4mm]}, thick}
    ]
    \node[box] (yaml) {YAML};
    \node[box, right=10pt of yaml] (po) {ParameterOrchestrator};
    \node[box, below=10pt of po] (ir)
      {Stream / VoiceManager / Controller$\times$4};
    \node[lbl, below=2pt of ir] (irlbl)
      {\textit{IR}: \texttt{List[List[Grain]]}};
    \node[box, below=18pt of irlbl] (renderer)
      {AudioRenderer (ABC)};
    \node[box, below left=10pt and -8pt of renderer] (np)
      {NumpyRenderer};
    \node[box, below=10pt of renderer] (cs)
      {CsoundRenderer};
    \node[box, below right=10pt and -10pt of renderer] (rp)
      {ReaperWriter};
    \node[lbl, below=2pt of np] (aif1) {\texttt{.aif}};
    \node[lbl, below=2pt of cs]  (aif2) {\texttt{.aif}};
    \node[lbl, below=2pt of rp]  (rpp)  {\texttt{.rpp}};
    \node[box, right=14pt of cs, align=center] (cache)
      {Stream\\CacheManager\\\fontsize{6}{7}\selectfont SHA-256};

    \draw[arr] (yaml) -- (po);
    \draw[arr] (po)   -- (ir);
    \draw[arr] (ir)   -- (irlbl);
    \draw[arr] (irlbl)-- (renderer);
    \draw[arr] (renderer) -- (np);
    \draw[arr] (renderer) -- (cs);
    \draw[arr] (renderer) -- (rp);
    \draw[arr] (np) -- (aif1);
    \draw[arr] (cs) -- (aif2);
    \draw[arr] (rp) -- (rpp);
    \draw[arr, dashed] (cache.west) -- (cs.east);
    \draw[arr, dashed] (cache.north west) to[bend left=15] (np.east);
  \end{tikzpicture}
  \caption{Pipeline del sistema. La specifica YAML è tradotta in IR
    Python dal \texttt{ParameterOrchestrator}; il layer di rendering
    è disaccoppiato dietro un'astrazione \texttt{AudioRenderer}; la
    cache SHA-256 opera per-stream sui renderer audio.}
  \label{fig:arch}
\end{figure}

\subsection{YAML come IR di intenzioni parametriche}

YAML è il \emph{syntax host}: lo schema semantico --- parametri,
inviluppi, range di variazione, distribuzioni, gate probabilistici,
strategie multi-voce --- è il \emph{domain language}. Il file YAML
non è una partitura nel senso che Truax critica
(\emph{«deterministic output [\dots] score or sequencer driven»},
\cite{Truax1988}, p.~19): è una specifica di intenzioni
parametriche tale che due esecuzioni dello stesso file producono,
con \texttt{dephase} attivo, output diversi. Un singolo stream è
descritto in forma compatta:

\begin{lstlisting}[language=,basicstyle=\ttfamily\fontsize{7.5}{9}\selectfont]
- stream_id: clouds
  onset: 0.0
  duration: 12.0
  sample: voce.aif
  grain_duration: {value: 0.08, range: 0.04}
  fill_factor: {value: 1.0, env: [0.3, 1.5, 0.5]}
  distribution: {value: 0.7}
  dephase: {grain_duration: 0.6}
\end{lstlisting}

Il programma di un'interfaccia dichiarativa ad alto livello per la
sintesi granulare è enunciato in ambito CIM da Di~Scipio e
Tisato~\cite{DiScipioTisato1993cim}: \emph{«a single rule may
instantiate multiple operations in the realization of an entire
process. This higher-level approach would represent a "step towards
the abstract"»} (p.~165). Roads
formula lo stesso programma articolato in tre punti distinti in
\emph{Microsound}~\cite{Roads2001} (cap.~5 p.~185, p.~234): un
linguaggio descrittivo musicale tradotto in parametri di
particelle, una specifica ad alto livello del montage, un controllo
mediante envelope, preset, automation. Vaggione~\cite{Vaggione1996}
ne articola la meccanica fine come \emph{«déclaration d'un attribut
particulier d'une entité quelconque [\dots] généralisé à toutes les
instances successives»} (p.~2): il valore YAML è la
\emph{déclaration}, l'inviluppo è la \emph{généralisation},
l'editing iterativo è l'\emph{imbrication} fra scrittura diretta e
trattamento algoritmico. Lo stesso programma è ripreso nella
tradizione CIM contemporanea da Markidis~\cite{Markidis2024cim}
(p.~48) come strumento di sustainability:
\emph{«by making the DSP score available [\dots] in a high-level
language not dependent on any specific implementation [\dots] the
composer abstracts the piece from its specific implementation
network»}.

Il \texttt{ParameterOrchestrator} traduce ogni chiave YAML in un
oggetto \texttt{Parameter} composto da una traiettoria centrale
\(c(t)\) (inviluppo time-varying), un range di deviazione
\(s\) (\texttt{mod\_range}) e una strategia di campionamento
(uniforme o gaussiana). Il modello di controllo è la
\emph{tendency mask} di Truax~\cite{Truax1988}: per ogni grano si
campiona \(v_n \sim \mathcal{D}(c(t_n), s)\) in modo indipendente
dal grano precedente, senza memoria di stato. Un secondo livello
--- il \texttt{ProbabilityGate} associato alla chiave
\texttt{dephase} --- decide \emph{se} la deviazione si applichi al
grano corrente, secondo una probabilità costante o time-varying.
Gate e range sono ortogonali: la \emph{déclaration} fissa il
valore, l'inviluppo lo generalizza nel tempo, il gate sceglie dove
la generalizzazione lascia spazio all'i.i.d.\ del grano.

Una conseguenza diretta è la disponibilità nel DSL di parametri
mutuamente esclusivi ma percettivamente complementari. \emph{Fill
factor} e \emph{density} controllano entrambi l'inter-onset time:
\texttt{density} è espresso in grani/secondo (controllo numerico
diretto), \texttt{fill\_factor} esprime la saturazione temporale
percepita ed è invariante alla durata del grano (\texttt{density}~$=
\texttt{fill\_factor}/\texttt{grain\_duration}$). Quale dei due
corrisponda all'intenzione compositiva in un contesto specifico è
una domanda empirica che solo il loop lungo --- partitura grafica e
ascolto --- permette di rispondere: la coppia è operativizzazione
diretta della tabella dei \emph{psychoacoustic correlates}
(Tab.~1 in~\cite{Truax1988}), non duplicazione.

Il sistema multi-voce è composto secondo la stessa logica
dichiarativa: $N$ voci sono differenziate su quattro assi
ortogonali (pitch, onset, pointer, pan) tramite strategie
selezionate da una singola chiave YAML per asse (linear,
geometric, stochastic, range, chord, spectral). Il
\emph{harmonization scheme} a $F$ voci di Truax~\cite{Truax1994}
viene generalizzato lungo quattro dimensioni indipendenti; la
\emph{décorrélation microtemporelle} di Vaggione~\cite{Vaggione2002}
emerge come effetto strutturale del layering, non come
post-processing.

\subsection{Renderer disaccoppiati, cache per-stream, export DAW}

Il layer di rendering implementa il pattern Open-Closed:
un'interfaccia astratta \texttt{AudioRenderer} è realizzata da tre
classi concrete. Il renderer nativo è
\texttt{NumpyAudioRenderer} (overlap-add Python in-memory, nessuna
dipendenza esterna); \texttt{CsoundRenderer} è un renderer
alternativo che invoca \texttt{csound} come sotto-processo a
partire da uno score testuale, fornito per chi opera nella
tradizione Roads/Csound; \texttt{ReaperProjectWriter} esporta un
progetto DAW senza renderizzare audio. La stessa specifica YAML è
servita da renderer diversi senza modifica --- precondizione
tecnica del programma di separazione fra specifica e implementazione
formulato in~\cite{Markidis2024cim}.

Due modalità di rendering sono disponibili. In modalità MIX tutti
gli stream confluiscono in un file unico con onset assoluti. In
modalità STEMS ogni stream è renderizzato come file indipendente,
con onset relativo a~0: è in questa modalità che si attiva la cache
incrementale. \texttt{StreamCacheManager} calcola il fingerprint
SHA-256 del dict YAML serializzato per stream (con
\texttt{sort\_keys=True} per indipendenza dall'ordine delle chiavi)
e considera \emph{dirty} uno stream se il fingerprint è cambiato,
se è nuovo, o se il file audio risulta assente su disco. Solo gli
stream dirty vengono ri-renderizzati; gli stream rinominati o
rimossi sono ripuliti via garbage collection del manifest. Su brani
con decine di stream e renderer Csound il tempo di build completa è
dell'ordine dei minuti: la cache non è un'ottimizzazione opzionale
ma la condizione tecnica perché il ciclo
\emph{modifica-un-parametro $\to$ riascolta} resti praticabile. Una
logica complementare di solo/mute per stream, dichiarata nel YAML
stesso, permette di isolare un singolo strato senza alterare la
struttura del file --- secondo meccanismo dello stesso ciclo.

L'export DAW chiude il workflow. \texttt{ReaperProjectWriter}
emette un file \texttt{.rpp} con una traccia per stream, posizione
e durata coerenti con il YAML, riferimento al file audio
corrispondente (al singolo stem in modalità STEMS, al mix unico in
modalità MIX). La struttura temporale dichiarata nel YAML transita
direttamente in DAW, senza ricostruzione manuale. Roads, descrivendo
la propria pratica di higher-order granulation assemblata in
Pro~Tools, lo formula così: \emph{«detached from real-time
constraints, ideas can be tested, edited, submixed, or deleted at
will»}~\cite{Roads2012}~(p.~8). Il sistema descritto istituzionalizza
questa pratica come architettura primaria, non come pipeline
manuale a posteriori.

\subsection{Language Server come scaffolding del loop lungo}

Lo schema semantico dei parametri --- nomi, tipi, range, gruppi
mutuamente esclusivi, valori ammessi per le strategie ---
non è un artefatto interno al codice di rendering ma è esposto come
JSON Schema e servito da un Language Server dedicato (denominato
\emph{[anonymous]-ls} nel repository). Il server è implementato in
Python sopra il protocollo Language Server Protocol via stdio, è
quindi integrabile con qualsiasi client compatibile; l'ambiente di
sviluppo di riferimento è VSCode tramite estensione dedicata. Tre
affordances sono offerte mentre il compositore scrive il YAML:
validazione inline (errori di tipo, range fuori bounds, mutua
esclusività violata, riferimenti a sample inesistenti); hover docs
sulle chiavi (descrizione del parametro, dominio, default);
completion contestuale sui valori (per esempio, lista degli accordi
disponibili sotto \texttt{chord\_pitch\_strategy}, finestre
registrate sotto \texttt{window}).

Il ruolo del Language Server nell'argomento del paper è preciso. Il
loop lungo nella sua forma minima --- specifica, generazione,
ascolto, riflessione, riscrittura --- ha un primo anello che
precede il rendering: la \emph{scrittura} stessa della specifica.
Senza scaffolding, ogni errore di sintassi o di dominio è scoperto
solo a rendering tentato, allungando ogni iterazione a tempi di
ordine non riducibile. Il LSP è il sito tecnico in cui
l'\emph{imbrication} écriture$\leftrightarrow$algorithme di
Vaggione~\cite{Vaggione1996} si materializza prima ancora che il
motore venga invocato: trasforma il YAML da formato testuale opaco
in interfaccia di scrittura compositiva con feedback ai più piccoli
incrementi della specifica.

% ================================================================
\section{La partitura grafica come strumento di retroazione}

% [DA SCRIVERE -- W3]
% Asse Y = posizione nel buffer; encoding visivo;
% uso nel loop lungo; confronto con precursori (Truax 1988 Fig. 4,
% Roads 1978/1988 polygon, Roads 2006 Ynez, EC2 Scan Display).
% Figura 2 (principale): screenshot partitura su brano reale.

% ================================================================
\section{Caso compositivo}

% [DA SCRIVERE -- W3, richiede brano + partitura]

% ================================================================
\section{Conclusioni}

% [DA SCRIVERE -- W4]

% ----------------------------------------------------------------
\section*{Riferimenti}
\bibliographystyle{plain}
\bibliography{refs}

\end{document}
